\chapter{Introdução à análise léxica}

\section{Introdução}

A análise léxica é a primeira fase do processo de compilação,
responsável por ler o código-fonte como uma sequência de caracteres e
agrupá-los em tokens significativos - palavras reservadas,
identificadores, operadores, números e símbolos - enquanto remove
elementos irrelevantes como espaços em branco e comentários. Sua
importância está em transformar o texto bruto do programa em uma
representação estruturada que simplifica as fases subsequentes
(análise sintática e semântica), além de detectar erros léxicos
precocemente, como caracteres inválidos ou identificadores mal
formados, funcionando como a ``porta de entrada'' que organiza e
valida o string de entrada antes das demais etapas de um compilador /
interpretador.

Damos o nome de \textbf{token} a um componente indivisível
da sintaxe de uma linguagem. De maneira simples, o conjunto de
tokens de uma linguagem pode ser entendido como o seu alfabeto.

Antes de apresentarmos a primeira técnica de análise léxica, vamos
descrever a primeira linguagem que iremos utilizar durante o curso.

\section{A linguagem Exp}

Durante o curso de compiladores vamos implementar interpretadores e
compiladores para diversas linguagens de programação. A primeira
destas linguagens será chamada de \textbf{Exp} e está é definida pela
seguinte gramática livre de contexto.

\[
  \begin{array}{lclclclcl}
    E & \to & n & | & E + E & | & E * E & | & (E)
  \end{array}
\]

De maneira simples, a linguagem \textbf{Exp} é formada por expressões
envolvendo adição, multiplicação, parêntesis e constantes numéricas
(representadas por \textbf{n}).

O conjunto de tokens desta linguagem é dado por \{n, +, *, (, )\}. A
próxima seção vamos apresentar a implementação de um analisador
léxico ad-hoc para esta linguagem.

\section{Analisadores léxicos ad-hoc}

Um analisador léxico ad-hoc é uma implementação manual e simplificada
de um analisador léxico, construída especificamente para um problema
ou linguagem particular, sem utilizar geradores automáticos como o
Alex. Diferentemente de analisadores baseados em expressões regulares
e autômatos (tema de um próximo capítulo), o analisador ad-hoc é
tipicamente implementado através de código que processa caracteres
sequencialmente, utilizando estruturas como casamento de padrão para
identificar tokens.

O primeiro passo de nossa implementação é a definição de tipos para
representação de tokens.

\begin{minted}{haskell}
type Line = Int
type Column = Int

data Token
  = Token {
      pos :: (Line, Column)
    , lexeme :: Lexeme
    }

data Lexeme
  = TNumber Int
  | TLParen
  | TRParen
  | TPlus
  | TTimes
  | TEOF
\end{minted}

Os sinônimos \haskell{Line} e
\haskell{Column} representam a informação de linha e
coluna inicial de um token. Por sua vez, tipo
\haskell{Token} é definido como um registro contendo
dois campos: um para sua posição inicial e outro que identifica o seu
\textbf{lexema}. O tipo \haskell{Lexeme} representa
os diferentes tokens da linguagem com seus construtores:
\haskell{TNumber} representa números,
\haskell{TLParen} o abre parêntesis,
\haskell{TRParen} o fecha parêntesis,
\haskell{TPlus} o operador de adição,
\haskell{TTimes} o operador de multiplicação e
\haskell{TEOF} o final da entrada.

O próximo passo para definirmos o analisador léxico para a linguagem
\textbf{EXP} é determinar o tipo desta função. A função deverá
receber uma string de entrada e retornar uma lista de tokens. Um
possível tipo para essa função seria:

\begin{minted}{haskell}
lexer :: String -> [Token]
\end{minted}

Porém, como lidar com situações de erro? Por exemplo, qual seria o
resultado se a string possuir caracteres inválidos (por exemplo,
\textbf{a}), que não fazem parte dos possíveis tokens da linguagem? O
ideal é que o analisador retorne uma mensagem de erro apropriada.
Para isso, vamos modificar o tipo de \haskell{lexer} para:

\begin{minted}{haskell}
lexer :: String -> Either String [Token]
\end{minted}

para permitir a emissão de mensagens de erro.

A ideia do analisador léxico é percorrer a string de entrada enquanto
produz uma lista contendo os tokens encontrados até o momento.
Durante o caminhamento da entrada, vamos atualizar o \emph{estado} do
analisador léxico a cada caractere lido. O estado é definido pelo
seguinte sinônimo:

\begin{minted}{haskell}
type State = (Line, Column, String, [Token])
\end{minted}

O tipo \haskell{State} é uma quádrupla formada pela
linha atual, coluna atual, string de dígitos consecutivos encontrados
até o momento e a lista de tokens já produzidos. A cada novo
caractere processado, devemos atualizar o estado de forma apropriada.
Para isso, utilizamos a função \haskell{transition},
apresentada a seguir.

\begin{minted}{haskell}
transition :: State -> Char -> Either String State
transition state@(l, col, t, ts) c
  | c == '\n' = mkDigits state c
  | isSpace c = mkDigits state c
  | c == '+' = Right (l, col + 1, "", mkToken state (Token TPlus (l,col)) ++ ts)
  | c == '*' = Right (l, col + 1, "", mkToken state (Token TMul (l,col)) ++ ts)
  | c == '(' = Right (l, col + 1, "", mkToken state (Token TLParen (l,col)) ++ ts)
  | c == ')' = Right (l, col + 1, "", mkToken state (Token TRParen (l,col)) ++ ts)
  | isDigit c = Right (l, col + 1, c : t, ts)
  | otherwise = unexpectedCharError l col c
\end{minted}

Quando encontramos parêntesis, operadores ou dígitos simplesmente
atualizamos o estado criando um token ou adicionando o caractere de
dígito na sequência de dígitos encontrados até o momento atual. Caso
o caractere atual seja um espaço ou uma quebra de linha, tentamos
construir um token de dígitos consecutivos e atualizamos a respectiva
posição de linha / coluna. A função
\haskell{mkDigits} realiza essa tarefa. Primeiro,
verificamos se a string de dígitos encontrados é vazia, nesta
situação retornamos o estado com a informação de linha / coluna
devidamente modificada. Caso a lista de dígitos seja formada apenas
por dígitos, criamos um novo token de dígitos e atualizamos o estado
apropriadamente. Finalmente, caso nenhuma destas situações se
aplique, temos um erro.

\begin{minted}{haskell}
mkDigits :: State -> Char -> Either String State
mkDigits state@(l, col, s, ts) c
  | null s = let l' = if c == '\n' then l + 1 else l
                 col' = if isSpace c then col + 1 else col
             in (l', col', s, ts)
  | all isDigit s = let t = Token (TNumber (VInt (read $ reverse s))) (l,col)
                        l' = if c == '\n' then l + 1 else l
                        col' = if c /= '\n' && isSpace c then col + 1 else col
                    in Right (l', col', "", t : ts)
  | otherwise = unexpectedCharError l col c
\end{minted}

De posse da função de transição entre estados do analisador léxico,
podemos definir o analisador léxico como:

\begin{minted}{haskell}
lexer :: String -> Either String [Token]
lexer = finish . foldl step base
  where
    base = Right (1, 1, "", [])

    finish = either Left (Right . extract)

    step ac@(Left _) _ = ac
    step (Right state) c = transition state c

    extract (l, col, s, ts)
      | null s = reverse ts
      | otherwise = let n = read (reverse s)
                        t = Token (TNumber (VInt n))
                                  (l, col - (length s))
                    in reverse (t : ts)
\end{minted}

O valor \haskell{base} é definido como a tupla
\haskell{(1, 1, "", [])} como valor do estado inicial
do analisador léxico formado por valores iniciais de linha, coluna e
as listas de tokens e dígitos iniciais é vazia.

A função \haskell{extract} é aplicada ao estado final
do analisador léxico para tentar construir um token de dígitos a
partir da sequência de dígitos consecutivos encontrados. Finalmente,
a função \haskell{finish} realiza o casamento de
padrão no resultado do analisar léxico: caso o resultado seja um
erro, este é propagado e caso seja o estado final do analisador,
aplicamos a função \haskell{extract} para obter a
lista final de tokens da entrada.

\section{Analisadores léxicos ad-hoc: vantagens e
desvantagens}

Analisadores léxicos ad-hoc possuem a vantagem de possuir uma
implementação simples para linguagens pequenas e maior controle sobre
erros encontrados. Porém, essa técnica possui diversas desvantagens.
A primeira é que tais analisadores são difíceis de manter e extender,
além de não serem eficientes para linguagens com uma estrutura léxica
complexa. Por exemplo, como estender a implementação do analisador
ad-hoc para prover suporte a identificadores? Como dar suporte a
comentários de linha ou mesmo de blocos?

Dessa forma, esta técnica é aplicável apenas em situações em que a
simplicidade da implementação supera as vantagens de um analisadores
léxicos produzidos automaticamente a partir de expressões regulares.

\section{Conclusão}

Neste capítulo apresentamos uma introdução à análise léxica e
descrevemos uma implementação de um analisador léxico ad-hoc para uma
linguagem de programação simples. Concluímos apresentando argumentos
a favor e contra o uso deste tipo de analisadores léxicos. Nos
próximos capítulos veremos como analisadores léxios podem ser
produzidos automaticamente a partir da descrição de tokens como
expressões regulares.
